\section*{Resumen}
\addcontentsline{toc}{section}{Resumen}

La utilización de descripciones de alto nivel para el prototipado rápido de algoritmos es una práctica recurrente tanto en la industria como en la academia. Esto se debe a que la abstracción y las herramientas disponibles facilitan el desarrollo, la depuración y la validación. Sin embargo, la mayoría de los prototipos debe implementarse en sistemas físicos embebidos que no cuentan con recursos para ejecutar directamente esas descripciones. En consecuencia, disponer de una herramienta que reduzca la brecha entre el modelado de alto nivel y la implementación embebida mejora los tiempos de desarrollo y la productividad, y además disminuye el error humano al sistematizar el proceso.

Esta memoria se enfoca en aportar a proyectos del ámbito académico. Para ello, explora la utilidad práctica de HDL Coder y HDL Verifier, herramientas de MATLAB para la generación automática de HDL desde scripts y diagramas de bloques dentro de un flujo de \textit{High-Level Synthesis} (HLS). En particular, se propone y valida una metodología sistemática para usar e implementar el código generado por HDL Coder en distintos casos de estudio ejecutados en FPGA sobre diferentes plataformas y entornos de verificación. En uno de los casos, relevante para proyectos en curso en el Departamento de Electrónica, se evalúan tiempos de ejecución y resultados de verificación en cada plataforma. Los códigos fuente y la información necesaria para reproducir y extender los experimentos están disponibles en un repositorio de código.

A partir de los resultados experimentales, se concluye que HDL Coder, apoyado por HDL Verifier, es una herramienta útil para generar código HDL optimizado para FPGA y mantener equivalencia funcional con las descripciones originales en MATLAB.
